\chapter{Semirings}

This chapter is dedicated to the study of semirings, which are algebraic structures that generalize rings by 
allowing for the absence of additive inverses. While semirings have been studied as early as the 1950s, 
the formal foundations for the algebraic geometry of semirings was largely developed only recently in \cite{borgersemi}.
The algebraic geometry for semirings presents many new challenges. For example, the notion of closed subschemes
is more subtle (see Remark~\ref{rem:badclosed} and Example~\ref{exmp:badclosed}).

\subsubsection*{Notation and conventions} The relevant module-theoretic framework for semirings has been established in \cite{borgersemi}.  
Let $\bb{N}$ denote the set of natural numbers including $0$ endowed with the usual addition and multiplication operators thus making
it a semiring. Then $\mod{\bb{N}}$ is the category of commutative monoids, and $\alg{\bb{N}}$ is the category of semirings. The usual
abstract module theory endows for any semiring $A$ its category $\mod{A}$ of $A$-modules.
Note that $\mod{A}$ is a symmetric monoidal category with tensor product $\otimes_A$ and unit object $A$, 
and furthermore, for any two modules $M,N \in \mod{A}$, $\Hom_{\bb{N}}(M,N)$ is an internal hom for $(\mod{A},\otimes_{A})$.\\
\indent For a semiring $A$, a module $M \in \mod{A}$ is said to be \textit{free} if it is isomorphic to a direct sum of copies of $A$.
A module $M$ is said to be \textit{flat} if the functor $-\otimes_A M:\mod{A} \to \mod{A}$ preserves finite limits. Later, it will be 
also useful to refer to this notion of flatness as \textit{categorical flatness}.

\begin{example}
Here are some examples of semirings that will serve as coefficient rings for our semiring schemes:
\begin{itemize}
\item $\bb{R}_{+}$, the semiring of non-negative real numbers with the usual addition and multiplication.
\item The Boolean semiring $\bb{B}$ is defined as the set $\{0,1\}$ with addition and multiplication defined as follows:
\[0+0=0, \quad 0+1=1, \quad 1+1=1\]
and
\[0\cdot 0=0, \quad 0\cdot 1=0, \quad 1\cdot 1=1.\]
\end{itemize}
\end{example}
For a natural number $n \in \bb{N}$ greater than $1$, let $[n]=\{1,2,...,n\}$.

\section{Preliminary results and motivation}
The semirings $\bb{N}$ and $\bb{R}_{+}$ naturally encode a notion of \textit{positivity} that flat modules over them inherit. 
We make this precise in the following definitions and propositions.
\begin{definition}
Let $M \in \mod{\bb{N}}$ be a commutative monoid. $M$ is said to be:
\begin{enumerate}
\item[(i)] \textit{Cancellative} if for all $m,m',m'' \in M$, $m+m'=m+m''$ implies $m'=m''$.
\item[(ii)] \textit{Negative-free} if for all $m,m' \in M$, $m+m'=0$ implies $m=0$ and $m'=0$.
\item[(iii)] \textit{Integral} if it is both cancellative and negative-free. 
\end{enumerate} 
\end{definition}
\begin{proposition}\label{prop:flatpositivity}
Let $A$ be a semiring, and $M$ a flat $A$-module. If $A$ is cancellative (resp. negative-free, integral), then so is $M$.
\end{proposition}
\begin{proof}
A commutative monoid $N$ is negative free if and only if:
\[0=\text{eq}(+,0: N\oplus N \to N)\]
where $+$ is the sum-map and $0$ is the zero map. Similarly, $N$ is cancellative if and only if the map
\[\text{id} \oplus \mathrm{diag}: N \oplus N \to N \oplus N \oplus N\]
taking an element $n\oplus n' \mapsto n \oplus n' \oplus n'$ realizes an isomorphism
\[N\oplus N = \text{eq}(f,g: N \oplus N \oplus N \to N)\]
where $f,g$ are defined as follows. For $n,n',n'' \in N$:
\[f(n \oplus n' \oplus n'') = (n+n', n'')\]
\[g(n \oplus n' \oplus n'') = (n+n'',n')\]
All propositions follow by considering the relevant equalizer diagram for $A$, which, 
after applying $\otimes_A M$, remains an equalizer diagram.
\end{proof}
\begin{corollary}\label{cor:NRnotflat}
The semiring maps $\bb{N} \to \bb{Z}$ and $\bb{R}_{+} \to \bb{R}$ are not flat.
\end{corollary}
Integral commutative monoids admit a natural partial order defined as follows:
\begin{definition}
Let $M$ be a commutative monoid. For $m,m' \in M$, we say that $m \le m'$ if there exists a $m'' \in M$ such that
\[m + m'' = m'.\] 
For an integral commutative monoid we will denote this comparison relation by $\le_M$.
\end{definition}
\begin{proposition}\label{prop:poset}
Let $M$ be an integral commutative monoid. Then $\le_M$ has the following properties:
\begin{enumerate}
\item[(i)] (Reflexivity) For all $m \in M$, $m \le_M m$.
\item[(ii)] (Transitivity) For all $m,m',m'' \in M$, if $m \le_M m'$ and $m' \le_M m''$, then $m \le_M m''$.
\item[(iii)] (Antisymmetry) For all $m,m' \in M$, if $m \le_M m'$ and $m' \le_M m$, then $m=m'$.
\item[(iv)] (Compatibility with addition) For all $m_1,m_2,m_1',m_2' \in M$, if $m_1 \le_M m_1'$ and $m_2 \le_M m_2'$, 
then $m_1 + m_2 \le_M m_1' + m_2'$.
\item[(v)] (Compatibility with scalar multiplication) For all $a \in A$ and $m,m' \in M$, if $m \le_M m'$, then $a m \le_M a m'$.
\item[(vi)] For all $m,m',m'' \in M$, if $m + m'' \le_M m' + m''$, then $m \le_M m'$.
\end{enumerate}
Therefore, $(M, \le_M)$ defines a partially ordered commutative monoid.
\end{proposition}
\begin{proof}
Out of the properties (i)-(v), only (iii) requires proof. Assume that $m \le_M m'$ and $m' \le_M m$. 
Then there exists $m_1,m_2 \in M$ such that $m= m' + m_1$ and $m' = m + m_2$. Thus $m=m+m_1+m_2$. 
Therefore $m_1+m_2=0 \implies m_1=m_2=0$, as desired.\\
\indent For (vi), assume that $m + m'' \le_M m' + m''$. Then there exists $m''' \in M$ such that $m + m'' + m''' = m' + m''$. 
By cancellation, we have $m + m''' = m'$, as desired, so $m \le_M m'$.
\end{proof}
On the other hand, categorical flatness is strong enough to tell us something about the existence of roots of polynomials.
\begin{proposition}
Let $\bb{R}_{+} \to A_{+}$ be a finite, flat map of semirings such that $A_+$ is reduced. Then $A_+$ has an $\bb{R}$-point, i.e. 
a map $A_{+} \to \bb{R}$ in $\alg{\bb{R}_{+}}$. 
\end{proposition}
\begin{proof}
It suffices to show that the groupification $A=A_{+} \otimes_{\bb{R}_{+}} \bb{R}$ admits a map to $\bb{R}$. Since $A_+$ is flat over $\bb{R}_{+}$,
$A_+ \to A$ is injective, and since $A_{+}$ is also reduced and finite over $\bb{R}_{+}$, we see that $A$ is finite and reduced itself, and is therefore 
a finite product of finite field extensions of $\bb{R}$, hence isomorphic to a finite product of copies of $\bb{R}$ and $\bb{C}$. Thus, in order to construct a map $A_{+} \to \bb{R}$, 
it suffices to construct a map $A \to \bb{R}$, and for that, it suffices to show that $A$ has at least one factor isomorphic to $\bb{R}$.
Assume for the sake of contradiction that $A$ has no $\bb{R}$-factor. Then $A$ is a product of copies of $\bb{C}$.\\
\indent Let $A_+$ be generated as an $\bb{R}_{+}$-module by elements $x_1, \ldots, x_n \in A_{+}$. For any $a \in A_+$, multiplication by 
$a$ defines an $\bb{R}_{+}$-linear map
\[m_a: A_{+} \to A_{+}.\]
Hence, we may view $m_a$ as an $n \times n$ matrix with entries in $\bb{R}_{+}$. Let 
\[S = \{ s = (s_1,...,s_n) \in \bb{R}_+^{n} \mid |s|_1 =\sum_{i \in [n]} s_i = 1\}\]
be the $(n-1)$-simplex, and consider the map $f: S \to S$ defined by
\[f(s) = \frac{m_a(s)}{|m_a(s)|_1}, \; \forall s \in S.\]
By the Brouwer fixed point theorem,
there exists a fixed point $r \in S$ such that $f(r) = r$,
and so we conclude that $m_a$ admits an eigenvector $r$ with a positive real eigenvalue $\lambda = |m_a(r)|_1$ 
(one could also appeal to the Perron-Frobenius theorem).\\
\indent Since $A$ is a product of copies of $\bb{C}$, there exists an element $a \in A$ whose eigenvalues are all non-real complex numbers.
However, by the above argument, $a$ must have a positive real eigenvalue, yielding a contradiction. Hence, $A$ must have at least one factor isomorphic to $\bb{R}$.
\end{proof}
\begin{corollary}
Let $A$ be a finite reduced $\bb{R}$-algebra. If $A$ has a positive model, i.e. there exists a finite, reduced and flat $\bb{R}_{+}$-algebra $A_{+}$, 
such that $A = A_{+} \otimes_{\bb{R}_{+}} \bb{R}$, then $A$ has at least one $\bb{R}$-point.
\end{corollary}
We therefore arrive at the following question 
that is closely related to the following version of Hilbert's Nullstellensatz in real algebraic geometry.
\begin{question}\label{mainquestion}
Let $A$ be a finitely presented flat $\bb{R}_{+}$-algebra. 
If $A$ is non-trivial, does it necessarily have an $\bb{R}$-point (or even better, an $\bb{R}_{+}$-point)?\\
\indent Equivalently, let 
\[\bb{R}^n_{+}:=\bb{R}_{+}[a_1,b_1,...,a_n,b_n]/f_n \sim g_n\]
where $f_n=1+\sum_{i=1}^n a_i^2 +b_i^2$ and $g_n=\sum_{i=1}^n 2a_i b_i$. By definition, $\bb{R}^n_{+}$ is the initial 
$\bb{R}_{+}[a_1,b_1,...,a_n,b_n]$-algebra $A$ such that the $f_n$ and $g_n$ are equal in $A$.\\
\indent Is it true that for every $n \ge 1$, there are no $\bb{R}_{+}$-algebra maps
\[\bb{R}^n_{+} \to A?\]
Let us explain the equivalence of the two questions. The existence of an $\bb{R}$-point of $A$ is equivalent to asking
for an $\bb{R}$-point of $A' = A \otimes_{\bb{R}_{+}} \bb{R}$. Picking a surjection $\bb{R}[x_1,...,x_n] \twoheadrightarrow A'$
with kernel $I$, $\spec{A'}$ does not admit $\bb{R}$-point if and only if $\text{rad}_{\bb{R}}(I) \neq \bb{R}[x_1,...,x_n]$, by the Real Nullstellensatz (see Theorem~\ref{realnull}). 
Therefore, $\spec{A'}$ admits $\bb{R}$-point if and only if $1 \notin \text{rad}_{\bb{R}}(I)$. Unpacking the definitions, 
and writing every element in $A'$ as the difference of two elements of $A$, we readily obtain the equivalence.
\end{question}
The equivalence mentioned above is due to the following version of Hilbert's Nullstellensatz in real algebraic geometry:
\begin{theorem}[Real Nullstellensatz]\label{realnull}
For an ideal $I \subseteq \bb{R}[x_1,...,x_n]$, let
\[V_{\bb{R}}(I) = \{x \in \bb{R}^n \mid f(x)=0 \text{ for all } f \in I\}\]
be its real variety, and 
\[\rad_{\bb{R}}(I)=\{f \in \bb{R}[x_1,...,x_n]:\; \exists m \in \bb{N}, f_1,...,f_d \in \bb{R}[x_1,...,x_n] :\;  -f^{2m}- (f_1^2 + ... + f_d^2) \in I\}.\]
Then $\rad_{\bb{R}}(I)$ is precisely the set of polynomials in $\bb{R}[x_1,...,x_n]$ that vanish on $V_{\bb{R}}(I)$. \cite[Chapter XI]{lang_algebra_2002}
\end{theorem}
Motivated by Question \ref{mainquestion}, we set-out the following goals:
\begin{enumerate}
\item[(a)] Develop a homotopical characterization of flatness for modules over semirings.
\item[(b)] Determine whether for a general flat finitely presented map of semirings $A \to B$, 
the map on spectra $\Spec{B} \to \Spec{A}$ is open.
\item[(c)] Determine when the equivalence relation generated by $f \sim g$ defines a flat map of semirings $\bb{R}_{+} \to \bb{R}_{+}[x_1,...,x_n]/(f \sim g)$.
\end{enumerate}
We will study these goals in the rest of the chapter.

\section{Homotopy theory of commutative monoids}
\indent Let $\Delta$ be the ordinary simplex category consisting of objects finite sets together with the usual set-maps.
For any category $\scr{C}$, we let $\scr{C}_{\Delta}$ denote the category consisting of objects being functors $\Delta^{\text{op}} \to \scr{C}$ and morphisms being 
natural transformations of functors; for any such functor $F: \Delta^{\text{op}} \to \scr{C}$, we call $F(n)$ the \textit{$n$-simplices} of $F$. For every object $C \in \scr{C}$, let $\underline{C} \in \scr{C}_{\Delta}$ be the constant pre-sheaf on $\Delta$ with value $C$; we will sometimes abuse notation and refer to such constant simplicial object $\underline{C}$ as simply $C$.\\
\indent We have a forgetful functor
\[U:\bb{N}\mathrm{-mod} \to \text{Set},\]
admitting a right adjoint
\[\bb{N}[-]:\mathrm{Set} \to \bb{N}\mathrm{-mod}\] 
assigning to every set $S \in \text{Set}$ the free commutative monoid $\bb{N}^{\oplus S}$. \\

\indent We also recall the Dold-Kan functor:
\[\DK:\text{Ch}(\mathrm{Ab}) \to \mathrm{Ab}_{\Delta}\]
As verified in \cite[III.2]{simpho}, the functor $\DK$ doesn't involve signs, so therefore descends to a well-defined functor
\[\DK:\text{Ch}(\bb{N}\mathrm{-mod}) \to \simpmod{\bb{N}}.\]
\begin{proposition}\label{prop:simpN}
Let $M,N \in \bb{N}\mathrm{-mod}_{\Delta}$ be simplicial commutative monoids, $S \in \mathrm{Set}_{\Delta}$ be a simplicial set, and $\Delta_{n}\in \mathrm{Set}_{\Delta}$ the simplicial set represented by $[n] \in \Delta$.
\begin{enumerate}
\item[(i)] The adjoint pair $(F \dashv U)$ extends to an adjoint pair:
\[\begin{tikzcd} \bb{N}\mathrm{-mod}_{\Delta} \arrow[r, bend right, "U"] & \text{Set}_{\Delta} \arrow[l, bend right, "F"]\end{tikzcd}\]
where $\bb{N}[S](n)=\bb{N}[S(n)]$.
\item[(ii)] The category $\bb{N}\mathrm{-mod}_{\Delta}$ may be endowed with a symmetric monoidal structure $\otimes_{\bb{N}}$, explicitly given by the following formula:
\[(M \otimes_{\bb{N}} N)(n) = M(n) \otimes_{\bb{N}} N(n).\]
There is an internal hom $\underline{\Hom}_{\bb{N}}$ given by the formula:
\[\underline{\Hom}_{\bb{N}}(M,N)(n) = \Hom_{\bb{N}\mathrm{-mod}_{\Delta}}(M \otimes_{\bb{N}}\bb{N}[\Delta_n],N).\]
 \end{enumerate}
 \end{proposition}
 \begin{proof}
 The proof of (i) is a standard fact for functor categories. For (ii), this follows from \cite[Tag017H]{sp}, after noting that $M \times \Delta_n$ is the same as $M \otimes_{\bb{N}}\bb{N}[\Delta_n]$ in our set-up.
 \end{proof}

We now wish to endow $\bb{N}\mathrm{-mod}_{\Delta}$ with a model structure in order to derive the tensor product $\otimes_{\bb{N}}$. 
Motived by the model structure on chain complexes that works so well, 
there is a natural such model structure, in hindsight perhaps suitably called the \textit{naive model structure}.
First we recall the \textit{homotopy-theoretic model structure} on $\text{Set}_{\Delta}$:
\begin{definition}[The homotopy-theoretic model structure on $\text{Set}_{\Delta}$]
A map $f: X \to Y$ in $\text{Set}_{\Delta}$ is a fibration if and only if it has the right lifting property with respect to all horn inclusions 
$\Lambda^n_k \to \Delta_n$ for all $n \ge 1$ and $0 \le k \le n$.\\
\indent The map $f: X \to Y$ in $\text{Set}_{\Delta}$ is a weak equivalence if and only if the induced map on geometric realizations
$|f|: |X| \to |Y|$ is a weak homotopy equivalence of topological spaces.\\
\indent The map $f:X \to Y$ is a cofibration if and only if it has the left lifting property 
with respect to all acyclic fibrations (i.e. fibrations that are weak equivalences).\\
\indent The map $f: X \to Y$ is an acyclic fibration if and only if it has the right lifting property with respect to all boundary inclusions
$\partial \Delta_n \to \Delta_n$ for all $n \ge 0$. \cite[I.11]{simpho}
\end{definition}
With this model structure on $\text{Set}_{\Delta}$, we may now define the \textit{naive model structure} on $\bb{N}\mathrm{-mod}_{\Delta}$.
\begin{definition}[The naive model structure]A map $f: X \to Y$ in $\bb{N}\mathrm{-mod}_{\Delta}$ is a fibration (resp. weak equivalence) 
if and only if $U(f)$ is a 
fibration (resp. weak equivalence) in $\text{Set}_{\Delta}$. 
Cofibrations are characterized to have the left lifting property with respect to all acyclic fibrations 
(i.e. fibrations that are weak equivalences).\\
\indent These three classes of 
morphisms define a model structure on $\bb{N}\mathrm{-mod}_{\Delta}$.
\end{definition}
This model structured has been considered many times in the literature, two notable references being \cite{kansemi} and \cite{rezksemi}.
Some immediate consequences of this definition are as follows.
\begin{proposition}\label{prop:hosimpN}
Let $X \in \simpmod{\bb{N}}$ be a simplicial commutative monoid.
\begin{enumerate}
\item[(i)] The adjoint functors $(F \dashv U)$ of Proposition~\ref{prop:simpN} form a Quillen adjunction for the naive model 
structure on $\simpmod{\bb{N}}$ and the homotopy-theoretic model structure on $\text{Set}_{\Delta}$. 
\item[(ii)] The adjoint functors $((X \otimes_{\bb{N}} -) \dashv \underline{\Hom}_{\bb{N}}(X,-)) : \simpmod{\bb{N}} \to \simpmod{\bb{N}}$ 
form a Quillen adjunction.
\item[(iii)] The category $\simpmod{\bb{N}}$ equipped with the naive model structure is a simplicial model category. 
In particular, if $Y \to Z$ is a cofibration in $\simpmod{\bb{N}}$, then so is $X \otimes_{\bb{N}} Y \to X \otimes_{\bb{N}} Z$.
\end{enumerate}
\end{proposition}
\begin{proof}
Both (i) and (ii) follow respectively from the fact that both $U$ and $\underline{\Hom}_{\bb{N}}(X,-)$ preserve both fibrations and acyclic fibrations. For $U$, 
this is immediate, and for $\underline{\Hom}_{\bb{N}}(X,-)$, this follows from the fact that it is a right adjoint, and the fact that fibrations and acyclic fibrations in $\simpmod{\bb{N}}$ are defined
via right lifting properties for the free $\bb{N}$-linear horn and boundary inclusions respectively (i.e. $\bb{N}[\Lambda^i_n] \hookrightarrow \bb{N}[\Delta_n]$ 
and $\bb{N}[\partial \Delta_n] \hookrightarrow \bb{N}[\Delta_n]$)\\
\indent For (iii), by \cite[Theorem II.5.1]{simpho} combined with \cite[Lemma II.6.1]{simpho}, it is enough to show that for any $X \in \simpmod{\bb{N}}$ 
we have a weak equivalence $X \to Q(X)$ with $Q(X) \in \simpmod{\bb{N}}$ and $U(Q(X))$ a Kan complex. This is provided to us by Kan's Extension Functor
that, to any simplicial set $S$, functorially associates a weak equivalence $S \mapsto \text{Ex}^{\infty}(S)$ where $\text{Ex}^{\infty}(S)$ a Kan complex.
\end{proof}
\begin{remark}
We may define the $\infty$-category $D_{\le 0}(\bb{N}\mathrm{-mod})$ to be the localization of simplicially-enriched category $\simpmod{\bb{N}}$ 
with respect to the weak equivalences defined by the naive model structure. Then $D_{\le 0}(\bb{N})$ is, up to 
equivalence of $\infty$-categories, the same as the \textit{animation} of $\bb{N}$-mod (see \cite[5.5.9.3]{htt}).
\end{remark}\label{rem:comptop}
With a model structure on $\simpmod{\bb{N}}$, we may proceed to calculate derived functors. In this paper, we will only be interested in doing so for $\otimes_{\bb{N}}$.
\begin{definition}
Let $N$ be an object of $\simpmod{\bb{N}}$. For any $M \in \simpmod{\bb{N}}$, we may take a cofibrant replacement $r: C(M) \to M$, and we define:
\[M \otimesl_{\bb{N}} N = C(M) \otimes_{\bb{N}} N,\]
as objects in $\Hot(\simpmod{\bb{N}})$.
\end{definition}
The naive model structure on $\simpmod{\bb{N}}$ is \textit{reasonable} in the following sense; by viewing $\bb{Z}\mathrm{-mod}_{\Delta} \hookrightarrow \bb{N}\mathrm{-mod}_{\Delta}$ 
as a full subcategory, then the induced model structure on $\simpmod{\bb{Z}}$ coincides with the usual model structure for chain complexes 
of abelian groups after applying Dold-Kan. However, there are quite a few conceptual differences between the two model-theoretic structures,
chief among them being the fact that objects in $\simpmod{\bb{N}}$ are rarely fibrant (i.e. Kan complexes) or even $\infty$-categories:
\begin{example}
\begin{enumerate}
\item[(a)] Objects of $\simpmod{\bb{N}}$ are rarely fibrant, or in other words, the underlying simplicial set is not a Kan complex.\\
\indent Indeed, while simple examples like $K(\bb{N},0) = \underline{\bb{N}}$, are Kan complexes, the next Eilenberg-Maclean space $K(\bb{N},1)$ is not \textemdash let us recall its definition. 
Consider the category $B\bb{N}$, which has one object and an endomorphism $1$ which acts freely; we define $K(\bb{N},1)=N(B\bb{N})$ where $N(B\bb{N})$ is its nerve. 
To show it is not a Kan complex, consider the map:
  \[\Lambda^2_{0} \to K(\bb{N},1)\]
  which takes the edge $[0,1] \to n$ and $[0,2] \to m$. If $K(\bb{N},1)$ were a Kan complex, then there would need to be a 
  number $h$ such that $n + h = m$, which is not solvable as soon as $n > m$.\\
  \indent Note that, while $K(\bb{N},1)$ is not a Kan complex, it is on the other hand an $\infty$-category, as is the nerve of any 
  ordinary category (see \cite[1.3]{keradon}).
\item[(b)] In fact, objects of $\simpmod{\bb{N}}$ are rarely $\infty$-categories. Indeed, we may construct a simplicial set 
modelling $K(\bb{N},2)$ which is not a quasi-category (see \cite{notquasi}). As $K(\bb{N},2)$ (which is the diagonal of the 
bisimplicial set $B^2 \bb{N}$) is modelled by an object in $\simpmod{\bb{N}}$ where each term is a finite free 
commutative monoid $\bb{N}^{\oplus r}$, 
from a homological standpoint, this example is essentially as nice as possible. 
We construct another example in Example~\ref{exmp:notquasi}.
\item[(c)] A surjection $X \twoheadrightarrow Y$ in $\simpmod{\bb{N}}$ need not be a fibration, even if it is termwise split on each degree.
\end{enumerate}
\end{example}
Motivated by trying to prove Question~\ref{mainquestion}, we wish to understand the structure of homotopy coequalisers in $\simpmod{\bb{N}}$. 
The structure of these equalizers is rather different from the case of abelian groups; in particular, we cannot show that they are always $1$-truncated. 
Calculating homotopy coequalizers amounts to understanding homotopy pushouts, which is greatly simplified in the case when the model category
is \textit{left proper} i.e. pushouts along cofibrations preserve weak equivalences. Fortunately, the naive model structure on $\simpmod{\bb{N}}$ is left proper. To show this, we could appeal to 
\cite[Theorem 8.1]{rezksemi}, which shows that it suffices to demonstrate that the functor 
\[ \simpmod{\bb{N}} \to \simpmod{\bb{N}}: \; \; M \mapsto M \oplus \underline{\bb{N}}\]
preserves weak equivalences, which can be done by hand. In particular, we note that for any $M \in \simpmod{\bb{N}}$ and $N \in \mod{\bb{N}}$,
\[M \oplus \underline{N} =\colim_{S' \subset U(N), |S'| < \infty} \prod_{S'} M = \hocolim_{S' \subset U(N), |S'| < \infty} \prod_{S'} M.\]
is a filtered colimit with cofibrant transition maps (i.e. monomorphisms in $\text{Set}_{\Delta}$) of products of copies of $M$, hence it preserves weak equivalences.
However, for completeness, we will give a more direct proof that follows the outline of Kan and Dwyer \cite{kansemi} 
(which unfortunately cannot be directly applied here as their categories do not result in commutative algebraic theories.)
\begin{proposition}
Let $M, M', N, N', K, K'$ be objects of $\simpmod{\bb{N}}$. 
\begin{enumerate}
\item[(i)]Let $f: M \to M'$ and $f': N \to N'$ be weak equivalences in $\simpmod{\bb{N}}$. 
Then the induced map
\[f \amalg f': M \oplus N \to M' \oplus N'\]
is also a weak equivalence.
\item[(ii)] Consider the pushout squares in $\simpmod{\bb{N}}$:
\[\begin{tikzcd}
M \arrow[r] \arrow[d] 
  & N \arrow[d] 
  && 
M' \arrow[r] \arrow[d] 
  & N' \arrow[d] \\
K \arrow[r] 
  & L
  \arrow[ul, phantom, "\ulcorner", very near start]
  &&
K' \arrow[r] 
  & L'
  \arrow[ul, phantom, "\ulcorner", very near start]
\end{tikzcd}\]
with weak equivalences $M \to M', N\to N', K\to K'$ making the diagrams commute, and moreover assume $M \to K$ and $M' \to K'$ are cofibrations.
Then the induced map $L \to L'$ is also a weak equivalence.
\item[(iii)]$\simpmod{\bb{N}}$ is left proper i.e. pushouts along cofibrations preserve weak equivalences.
\end{enumerate}
\end{proposition}
\begin{proof}
For (i), it suffices to prove the case that $M=M'$. We have already established the case when $M$ is discrete.
We may then proceed using the diagonal lemma for bisimplicial sets, as by definition, $M \oplus N = \mathrm{diag}(L_{\bullet, \bullet})$ where
\[L(i,j) = M(i) \oplus N(j).\]
The result then follows from the fact that for each fixed $i$, the map $M(i) \oplus N \to M(i) \oplus N'$ is a weak equivalence, so the diagonals
are weak equivalences.\\
\indent For (ii), we can follow \cite[Proposition 8.1]{kansemi} with the follow modifications. First, their free-product $*$ serves the same purpose as
our direct sum $\oplus$ in $\simpmod{\bb{N}}$. Their Lemma 2.7 is then replaced by our (i) and for their Proposition 8.2, we need
a characterization of cofibrations in $\simpmod{\bb{N}}$ similar to that of their 7.6. After noting that their free categories $F$ serve
the same functorial purpose of our $\bb{N}[-]$, then as $\bb{N}[\partial \Delta_n] \hookrightarrow \bb{N}[\Delta_n], \forall n \ge 0$ form a set of generating cofibrations,
their 7.6 also holds in our setting. The proof then goes through verbatim.\\
\indent Finally, (iii) follows from (ii) since any pushout diagram can be made into a cofibrant diagram which preserves weak equivalences.
\end{proof}
The following theorem describes our current understanding of the structure of homotopy coequalizers in $\simpmod{\bb{N}}$.
\begin{theorem}\label{thm:main0}
Let $f,g: M \to N$ be two maps in $\bb{N}\mathrm{-mod}$. There is a $Z(f,g) \in \simpmod{\bb{N}}$ such that $Z(f,g) \simeq \hocoeq(f,g)$ and:
\begin{enumerate}
\item[(i)] $Z(f,g)$ is a $(2,1)$-category that is weakly $1$-coskeletal and minimal in dimension $1$.
\footnote{In particular, $Z(f,g)$ as a simplicial set is isomorphic to the Duskin nerve of a $2$-category in which every 
$2$ morphism is an isomorphism, see \cite[4.8.1.5]{keradon}}
\item[(ii)] If $Z(f,g)$ is an $\infty$-category, then $Z(f,g)$ is further a $(1,1)$-category and $Z(f,g)=N(\Hot(Z(f,g)))$. For example, this is always true when $g=0$.
\item[(iii)] If $g=0$, $N$ is cancellative and $M$ is negative-free, then $Z(f,g)$ is discrete and $\hocoeq(f,g)=\coeq(f,g)$.
\end{enumerate}
\end{theorem}
\begin{remark}
It is worth comparing Theorem~\ref{thm:main0} with the usual structure in abelian groups. 
Indeed, the homotopy coequaliser of two maps $f,g: M \to N$ is quasi-isomorphic to the $\hocoker(f-g)$, 
which is concentrated in degrees $0$ and $1$. 
Our theorem recovers this fact. 
In particular, (i) says that $\hocoeq(f,g)=Z(f,g)$ is a $(2,1)$ category, and we know that $Z(f,g)$ is a Kan complex, so (ii) says that $Z(f,g)$ 
is in fact a $(1,1)$ category. Therefore, $Z(f,g)$ is a $1$-groupoid, and in particular, is $1$-truncated (see \cite[3.5]{keradon}).
\end{remark}
To prove Theorem~\ref{thm:main0}, we need to recall the following definition:
\begin{definition}\label{def:n1cat}
Let $n$ be a positive integer. A simplicial set $S$ is an $(n,1)$-category if it satisfies the following condition for every pair of integers $0 < i < m$: 
the restriction map 
\[\Hom_{\text{Set}_{\Delta}}(\Delta^m, S) \to \Hom_{\text{Set}_{\Delta}}(\Lambda_i^m, S)\]
is surjective, and if $m > n$, is also injective. A simplicial set $S$ is weakly $n$-coskeletal if the restriction map
\[\Hom_{\text{Set}_{\Delta}}(\Delta^m, S) \to \Hom_{\text{Set}_{\Delta}}(\partial \Delta^m, S)\]
is bijective for $m \ge n+2$ and an injection for $m=n+1$.
\end{definition}
\begin{proposition}\label{prop:11cat}
A simplicial set $S$ is isomorphic to the nerve of a category $\cal{C}_0$ if and only if it is a $(1,1)$-category. 
In particular, $S=N(\Hot(S))$.
\end{proposition}
\begin{proof}
See \cite[Proposition 1.3.4.1]{keradon}.
\end{proof}
\begin{proposition}\label{prop:coeqN}
Let $f,g: X \to Y$ be two morphisms in $\simpmod{\bb{N}}$, and let $Z(f,g)$ fit into a pushout diagram in $\simpmod{\bb{N}}$:
\[\begin{tikzcd}
X \amalg X \arrow[d, "\text{id}_X \otimes {\bb{N}[\iota]}"] \arrow[r, "(f\text{,}g)"] & Y \arrow[d]\\
X \otimes_{\bb{N}} {\bb{N}[\Delta^1]}\arrow[r] & Z(f,g)
\end{tikzcd}\]
where $\iota: \partial \Delta^0 \hookrightarrow \Delta^1$ is the canonical inclusion as a sub-simplicial set. Then $Z(f,g) \simeq \hocoeq(f,g:X \to Y)$.
\end{proposition}
\begin{proof}
Let $Z'=X \coprod^{\text{h}}_{\nabla,X \amalg X, (f,g)} Y$, note that $Z' \simeq \hocoeq(f,g:X \to Y)$. To show the claim, 
it suffices to show that $\text{id}_X \otimes \bb{N}[\iota]: X\amalg X \to X \otimes_{\bb{N}} \bb{N}[\Delta^1] $ is a cofibrant replacement for $\nabla: X \amalg X \to X$. This will follow from the following two claims:
\begin{enumerate}
\item[(i)] The canonical map $X \otimes_{\bb{N}} \bb{N}[\Delta^1] \to X$ induced by $\Delta^1 \to *$ is a homotopy equivalence.
\item[(ii)] The map $\text{id}_X \otimes \bb{N}[\iota]: X\amalg X \to X \otimes_{\bb{N}} \bb{N}[\Delta^1]$ is a cofibration.
\end{enumerate}
Note that (ii) follows from (i) and (iii) of Proposition~\ref{prop:hosimpN}, so it suffices to show (i). 
But this can be shown after applying $U$, so it suffices to prove a weak equivalence between the simplicial sets 
$\underline{\Hom}_{\text{Set}_{\Delta}}(\Delta^1, X)$ and $\underline{\Hom}_{\text{Set}_{\Delta}}(\Delta^0, X)$, which is clear as $\Delta^1 \to \Delta^0$ is a homotopy equivalence.
\end{proof}
\begin{example}\label{exmp:coeqN}
Consider the situation of Proposition~\ref{prop:coeqN} specialized to the case where $X=\underline{M}$ and $Y=\underline{N}$, where $M$ and $N$ are $\bb{N}$-modules. 
Then $Z(f,g)$, as a functor $\Delta^{op} \to \bb{N}$ (ignoring degeneracy maps), may be described as follows: $\forall (m_1,...,m_n) \in M^{\oplus n}, k \in N$
\begin{align*}
Z(f,g)([n]) &= M^{\oplus n} \oplus N, \; \; \forall [n] \in \ob(\Delta^{op}),\\
Z(f,g)(\partial_i)(m_1,...,m_n,k) &= (m_1,...,m_{i-1}, m_i + m_{i+1},...,m_n, k), \; \;  \forall 0 < i < n,\\
Z(f,g)(\partial_0) (m_1,...,m_n,k) &= (m_2,...,m_n, f(m_1)+k),\\
Z(f,g)(\partial_n) (m_1,...,m_n,k) &= (m_1,...,m_{n-1}, g(m_n)+k).
\end{align*}
$Z(f,g)$ has some nice features, as the next proposition shows.
\end{example}
\begin{proposition}\label{prop:weakskel}
Let $Z(f,g)$ be the simplicial set described in Example~\ref{exmp:coeqN}, i.e. $\forall (m_1,...,m_n) \in M^{\oplus n}, k \in N$
\begin{align*}
Z(f,g)([n]) &= M^{\oplus n} \oplus N, \; \; \forall [n] \in \ob(\Delta^{op}),\\
Z(f,g)(\partial_i)(m_1,...,m_n,k) &= (m_1,...,m_{i-1}, m_i + m_{i+1},...,m_n, k), \; \;  \forall 0 < i < n,\\
Z(f,g)(\partial_0) (m_1,...,m_n,k) &= (m_2,...,m_n, f(m_1)+k),\\
Z(f,g)(\partial_n) (m_1,...,m_n,k) &= (m_1,...,m_{n-1}, g(m_n)+k),
\end{align*}
where $M, N \in \bb{N}\mathrm{-mod}$ and $f,g: M \to N$ are morphisms in $\bb{N}\mathrm{-mod}$.\\
\indent Then for any $m \ge 3$ and $0 < i < m$, the restriction maps
\[a^m_i: \Hom_{\simpmod{\bb{N}}}(\Delta^m, Z(f,g)) \to \Hom_{\simpmod{\bb{N}}}(\Lambda_i^m, Z(f,g))\]
\[b^m: \Hom_{\simpmod{\bb{N}}}(\Delta^m, Z(f,g)) \to \Hom_{\simpmod{\bb{N}}}(\partial \Delta^m, Z(f,g))\]
are bijective, and injective for $m=2$.\\
\indent If $g=0$, then we also have bijectivity for $a^2_1$, so $Z(f,0) \simeq N(\Hot(Z(f,0)))$, 
and in particular, $Z(f,0)$ is an $\infty$-category. $\mathrm{Hot}(Z(f,0))$ can be described as follows:
\begin{itemize}
  \item $\ob(\Hot(Z(f,0))) = \{n| n'\in N\}$
  \item $\Mor(\Hot(Z(f,0)))(n,n') = \{m \in M | n'=n+f(m)\}$.
  \item The composition rule 
  \[\Mor(\Hot(Z(f,0)))(n',n'') \times \Mor(\Hot(Z(f,0)))(n,n') \to \Mor(\Hot(Z(f,0)))(n,n'')\]
  is described via addition on the commutative monoid $M$.
\end{itemize}
\end{proposition}
\begin{proof}
Let us first assume $m\ge 3$. We note that any $x=(m_1,...,m_m, n) \in Z(f,g)([m])$ is determined by 
\begin{align*}
\partial_0 x &= (m_2,...,m_m, f(m_1)+n)\\
\partial_n x &= (m_1,...,m_{m-1}, g(m_m)+n)\\
\partial_j x &= (m_1,...,m_{j-1}, m_j + m_{j+1},...,m_m, n)
\end{align*}
for any $0 < j < m$ and $j \neq i$. Thus, injectivity is clear, and for surjectivity, forming a candidate for $x$ given 
$\partial_j x$ for $j \neq i$ is straightforward by these equations, and the fact that it is indeed a lift to $Z(f,g)([m])$ of the boundaries
follows from the simplicial identities in both situations. \\
\indent Now, let us assume $m=2$. For injectivity, let $x=(m_1,m_2,n) \in M_{\bb{Z}}^{\oplus 2} \oplus N_{\bb{Z}}$, then:
\[\partial_0 x = (m_2,f(m_1)+n), \partial_2 x  = (m_1, n),\]
so indeed, $x$ is determined by $\partial_0 x, \partial_2 x$ and thus injectivity for $m=2$.\\
\indent If $g=0$, then the simplicial identities imply surjectivity for $a^2_1$. To be explicit, given $(m_1,n_1), (m_2,n_2)$ 
with $f(m_2)+n_2=n_1$ forming an element $\Lambda^2_1 \in Z(f,g)([1])$, then $x=(m_2,m_1,n_2)$ is an element in $Z(f,g)([2])$ with
$\delta_0x=(m_1,n_1)$ and $\delta_2x=(m_2,n_2)$.\\
\indent Therefore, $Z(f,0)$ is a $(1,1)$-category (see Definition~\ref{def:n1cat}), and thus $Z(f,0) \simeq N(\Hot(Z(f,0)))$ by Proposition~\ref{prop:11cat}. 
The description of $\Hot(Z(f,0))$ is then straightforward to verify.
\end{proof}
\begin{remark}\label{rem:mindim1}
We note that the simplicial set $Z(f,g)$ of Proposition~\ref{prop:weakskel} is also minimal in dimension $1$ (see \cite[Definition 4.7.6.4]{keradon}). 
This amounts to the following verification: given $x,x' \in Z(f,g)([1])$, if there exists simplices $y,y' \in Z(f,g)([2])$ such that
\[\partial_0 y =x, \partial_2 y = s(\partial_1 x), \; \partial_2 y' =x', \partial_0 y' = s(\partial_0 x'), \; \partial_1 y = \partial_1 y',\]
then $x=x'$. In our case, we can just directly verify:
\[x=\partial_1 y = \partial_1 y' = x'.\]
\end{remark}
\begin{example}\label{exmp:notquasi}
One may ask whether all simplicial sets $Z(f,g)$ constructed in Example~\ref{exmp:coeqN} are $\infty$-categories. 
We will show that this is not the case by consider $M=N=\bb{N}$ and $f,g: \bb{N} \to \bb{N}$ given by $f(n)=2n$ and $g(n)=n$ (note that $\bb{B}= \coeq(f,g: \bb{N}\to \bb{N})$).\\
\indent In particular, we remark that there exists a morphism $n \to n' \in Z(f,g)([0])=\bb{N}$ if and only if $n \le n' \le 2n$. 
So for example, there is an inner horn $\Lambda^1_2$ with edges $ 1 \to 2 \to 3$ in $Z(f,g)([1])$, but since there is no morphism from $1 \to 3$, 
there cannot be a $2$-simplex in $Z(f,g)([2])$ filling this horn.
\end{example}
\begin{proof}[Proof of Theorem~\ref{thm:main0}]
Using Proposition~\ref{prop:coeqN}, Example~\ref{exmp:coeqN} and Proposition~\ref{prop:weakskel}, 
we find a simplicial set $Z(f,g)$ which is a $(2,1)$-category that is weakly $1$-coskeletal 
such that $Z(f,g) \simeq \hocoeq(f,g:M \to N)$.
The second part of Proposition~\ref{prop:weakskel} implies that $Z(f,g)$ is an $\infty$-category when $g=0$, 
and indeed, whenever $Z(f,g)$ is an $\infty$-category, $Z(f,g)$ is a $(1,1)$-category 
by \cite[Proposition 4.8.1.7]{keradon} and Remark~\ref{rem:mindim1}, so by \cite[Example 4.8.1.3]{keradon}, $Z(f,g) \simeq N(\Hot(Z(f,g)))$.\\
\indent Finally, when $g=0$, $N$ is cancellative and $f$ is injective, then every two objects of $\Hot(Z(f,0))$ are 
connected by a unique morphism, and if $M$ is negative-free,
then there are no non-trivial automorphisms, so $\Hot(Z(f,0))$ is a poset. There is then a 
contracting homotopy taking $N(\Hot(Z(f,0)))$ to its set of connected components, and thus 
$\hocoeq(f,0)=\coeq(f,0)$.
\end{proof}
As an immediate application of Theorem~\ref{thm:main0}, we have the following:
\begin{theorem}\label{thm:ZN}
We have 
\[\bb{Z} \simeq \hocoeq(\mathrm{diag},0: \bb{N} \to \bb{N} \oplus \bb{N}).\]
Therefore:
\[\bb{Z} \otimesl_{\bb{N}} \bb{Z} \simeq \bb{Z}[0].\]
\end{theorem}
\begin{proof}
The conditions of Theorem~\ref{thm:main0} (iii) are satisfied, so we have:
\[\hocoeq(\mathrm{diag},0: \bb{N} \to \bb{N} \oplus \bb{N}) \simeq \coeq(\mathrm{diag},0: \bb{N} \to \bb{N} \oplus \bb{N}) \simeq \bb{Z},\]
Since the functor $-\otimesl_{\mathbb{N}} \mathbb{Z}$ commutes with homotopy coequalizers:
\[\bb{Z} \otimesl_{\bb{N}} \bb{Z} \simeq \hocoeq(\mathrm{diag},0: \bb{N} \to \bb{N} \oplus \bb{N}) \otimesl_{\bb{N}} \bb{Z}
\simeq \hocoeq(\mathrm{diag},0: \bb{Z} \to \bb{Z} \oplus \bb{Z}) \simeq \bb{Z}[0],\]
as desired.
\end{proof}
It is also interesting to relate the explicit construction of $Z(f,g)$ in Example~\ref{exmp:coeqN} to other potential constructions.
\begin{corollary}\label{cor:DKcoeq}
Let $f,g: X \to Y$ be two maps in $\simpmod{\bb{N}}$. Let $\Delta^{nd}_{\le 1}$ denote the subcategory of $\Delta_{\le 1}$ 
where we do not include the degeneracy map $[1] \to [0]$. 
Consider the canonical inclusion of categories $i:\Delta^{nd,op}_{\le 1} \to \Delta^{op}$, 
and let $F:\Delta^{nd,op}_{\le 1} \to \simpmod{\bb{N}}$ be the functor representing $f,g:X \to Y$.\\
\indent If $F'$ is the homotopy left Kan extension of $F$ along $i$, then $F' \simeq \hocoeq(f,g: X \to Y)$. 
Moreover, if $g=0$ and $X=\underline{X([0])}, Y=\underline{Y([0])}$ then $F'\simeq \DK(f: X[0] \to Y[0])$.
\end{corollary}
\begin{proof}
A defining property of homotopy left Kan extensions is that 
\[\hocolim_{\Delta^{nd}_{\le 1}} F = \hocolim_{\Delta^{op}} F',\]
so the first claim follows directly. For the second claim, by Proposition~\ref{prop:coeqN} and Example~\ref{exmp:coeqN}, 
we have an explicit description of $\hocoeq(f,g: X \to Y)$ using the simplicial set $Z(f,0)$ defined there, and in this case, 
one directly verifies that as simplicial sets, $Z(f,0)$ and $\DK(f: X[0]\to Y[0])$ are the same, resolving the second claim.
\end{proof}
The following theorem is rather surprising.
\begin{theorem}\label{thm:contractab}
Let $\iota_M: M \to M \otimes_{\bb{N}} \bb{Z}$ be the canonical map of a commutative monoid $M \in \mod{\bb{N}}$ 
to its groupification. Then $\hocoeq(\iota_M,0: M \to M \otimes_{\bb{N}} \bb{Z})$ is contractible.
\end{theorem}
\begin{proof}
By Theorem~\ref{thm:main0}, we know that $\hocoeq(\iota_M,0: M \to M \otimes_{\bb{N}} \bb{Z}) \simeq N(\scr{C})$ 
where $\scr{C}$ with objects $M \otimes_{\bb{N}} \bb{Z}$ and morphisms $n \mapsto n + i(m)$ for $m \in M$ (see Proposition~\ref{prop:weakskel}). 
We claim that $\scr{C}$ is a filtered category,
which will be enough to prove the desired contractibility of $N(\scr{C})$ (see \cite[Proposition 5.3.1.18]{htt}). We note the following trivial observation: 
for any $m \in M \otimes_{\bb{N}} \bb{Z}$, either $m$ or $-m \in \im(\iota_M)$.\\
\indent First, we show that for any two arrows $f,g: X \to Y \in \scr{C}$, there is a $Z \in \scr{C}$ and a map $h:Y \to Z$ such that $hf=hg$.\\
\indent The two maps $f,g$ may be represented as maps $n \to n + \iota_M(m)$ and $n \to n+\iota_M(m')$, 
where $n=X \in M \otimes_{\bb{N}} \bb{Z}$ and $m,m' \in M$, 
and we have the relation $n+\iota_M(m) = n + \iota_M(m')$ so $ \iota_M(m)  = \iota_M(m')$. Therefore, there exists a $k \in M$ 
such that $k + m' = k+m$. 
Take $Z = n + \iota_M(k + m) = n+\iota_M(k + m')$, and the map $Y \to Z$ given by $Y \to Y + \iota_M(k)$.\\
\indent Next we show that for any two elements $X, Y \in \scr{C}$, there is a third $Z \in \scr{C}$ and maps $X \to Z \leftarrow Y$.\\
\indent We represent $X, Y$ by two elements $n, n' \in M \otimes_{\bb{N}} \bb{Z}$. 
If $n$ and $n'$ are both in the image of $\iota_M$, so that $n = \iota_M(m)$ and $n' = \iota_M(m')$, 
then we take $Z = \iota_M(m + m')$. On the other hand, if $X$ is not in the image of $\iota_M$, then $-n=\iota_M(m)$, 
so there is an arrow $X \to 0$ given by $n \mapsto n + \iota_M(m)$. If $Y$ is in the image of $\iota_M$, then there is an arrow $X \to Y$ so we set $Z=Y$,
and otherwise, we set $Z=0$.\\
\indent These two facts combined are enough to show the claim.
\end{proof}
\begin{corollary}\label{cor:main}
For any $M \in \bb{N}\mathrm{-mod}$, we have $M \otimesl_{\bb{N}} \bb{Z} \simeq M \otimes_{\bb{N}} \bb{Z}[0]$.
\end{corollary}
\begin{proof}
By Theorem~\ref{thm:contractab} we have:
\[\hocoeq(\iota_M,0: M \to M \otimes_{\bb{N}} \bb{Z}) \simeq 0.\]
By Theorem~\ref{thm:ZN}, we have:
\begin{align*}
\hocoker(\iota_M: M \otimesl_{\bb{N}} \bb{Z} \to M \otimes_{\bb{N}} \bb{Z}) & \simeq \hocoeq(\iota_M,0: M \otimesl_{\bb{N}} \bb{Z} \to M \otimes_{\bb{N}} \bb{Z} \otimesl_{\bb{N}} \bb{Z})\\
&\simeq \hocoeq(\iota_M,0: M \to M \otimes_{\bb{N}} \bb{Z}) \otimesl_{\bb{N}} \bb{Z}\\
&\simeq 0.
\end{align*}
Thus, the map $\iota_M: M \otimesl_{\bb{N}} \bb{Z} \to M \otimes_{\bb{N}} \bb{Z}$ is a weak equivalence, as desired.
\end{proof}
\section{Flat topology for semiring schemes}
The previous section showed that homotopy theory is not enough to understand the algebraic structure of flatness. We thus now
turn to studying flatness in a more algebro-geometric manner. 
\begin{definition}[Zariski topology on semirings]
For a semiring $R$, an \textit{ideal} of $R$ will just be an $R$-submodule. We say an $R$-submodule $\ideal{p} \subset R$ is \textit{prime} if $R \setminus \ideal{p}$ is a multiplicatively closed subset. 
Furthermore, for any multiplicatively closed subset $S \subset R$, we may define the 
\textit{localization} semiring $S^{-1}R$ in the usual way: elements of $S^{-1}R$ are represented as `fractions' $r/a$ 
under the equivalence relation 
\[r/s \sim r'/t \Longleftrightarrow \exists q \in S \text{ such that } qtr = qsr'.\]
We may then define the semiring $R[1/a]$ for any $a \in R$ as the localization along the multiplicatively closed subset $\{1,a,a^2,...\} \subset R$.\\
\indent Let $\text{Spec}(R)$ denote the topological space whose points are prime ideals $\ideal{p} \subset \text{Spec}(R)$ equipped with the \textit{Zariski} topology, 
the topology generated by the subsets $D(a) = \{\ideal{p} \cap R \; | \; \ideal{p} \subset \text{Spec}(R[1/a])\}$.\\
\indent A semiring map $f: A \to B$ induces a map $\spec{f}: \spec{B} \to \spec{A}$ via the association $\ideal{q} \mapsto f^{-1}(\ideal{q})$.\\
\indent A semiring $A$ is \textit{local} if the following equivalent conditions hold:
\begin{enumerate}
\item[(a)] $\spec{A}$ has a unique closed point.
\item[(b)] $A$ has a unique maximal ideal.
\item[(c)] There is a proper ideal $\ideal{m} \subset A$ such that every element in $A \setminus \ideal{m}$ is invertible.
\end{enumerate}
One may check the equivalence of these conditions as in classical commutative algebra.
\end{definition}
\begin{proposition}\label{prop:spec}
Let $A\to B$ be a map of semirings. The following hold:
\begin{enumerate}
\item[(i)] Closed subsets of $\spec{A}$ correspond to ideals $I \subset A$ via the association $I \mapsto V(I) = \{\ideal{p} \; | \; I \subset \ideal{p}\}$.
\item[(ii)] $\spec{A}$ is a spectral space, and therefore has an associated constructible topology.
\item[(iii)] The map $\spec{B} \to \spec{A}$ induced by $A \to B$ is continuous with respect to the Zariski topology, and is spectral.
\end{enumerate}
\end{proposition}
\begin{proof}
The proofs of (i)-(iii) are similar to those in classical commutative algebra, see for example \cite[Tag 00DY]{sp}.
\end{proof}
On the other hand, we may define the \textit{fppf topology} on $\text{Spec}(R)$. 
\begin{definition}[The fppf topology on semirings]
Given an $R$-algebra $f:R \to S$, we say $S$ is a \textit{flat $R$-algebra} 
if it is flat as an $R$-module under the structure map $f$, and \textit{faithfully flat} if $S$ is additionally faithful i.e. 
for any map of $R$-modules $f: M \to N$, $f$ is an isomorphism if and only if the induced map 
$f \otimes_R S: M \otimes_R S \to N \otimes_R S$ is injective.\\
\indent An $R$-algebra $S$ is finitely presented if it can be written as a coequaliser of two $R$-algebra homomorphisms 
$f,g: R[y_1,...,y_r] \to R[x_1,...,x_d]$. Equivalently, the functor $\alg{R} \to \text{Set}$ taking an $R$-algebra $A$ to the set 
$\Hom_{\alg{R}}(S, A)$ commutes with filtered colimits.
\end{definition}
With some care, we can adapt arguments in classical commutative algebra. The main difference is that elements in a commutative monoid
no longer `communicate' with each other (via subtraction), so we must work with elements more directly. As a concrete example,
there is no single submodule that detects whether a map is injective $f: M \to N$, rather we must work with all the fibers $f^{-1}(n) \subset M$
for all $n \in N$.
\begin{proposition}\label{prop:flatsurjmax}
Let $f: A \to B$ be a flat map of semirings such that $\text{Spec}(f):\text{Spec}(B) \to \text{Spec}(A)$ is surjective on maximal ideals, then $f$ is faithfully flat.
\end{proposition}
\begin{proof}
We must show that if $\phi: M \to N$ is map of $A$-modules, then 
\[\phi \otimes_A B \text{ is an isomorphism} \implies \phi \text{ is an isomorphism}.\]
We shall first show that $\phi$ is injective.
\begin{claim}\label{claim:flatsurjmax1}
For any $A$-module $M$, the canonical map $M \to M \otimes_A B$ sending $m \mapsto m\otimes 1$ is injective.
\end{claim}
\begin{proof}[Proof of claim~\ref{claim:flatsurjmax1}]
First, we note that if $I \hookrightarrow A$ is an ideal, so a monomorphism of $A$-modules, 
then since $A \to B$ is flat, $I \otimes_A B = IB$.\\
\indent For any two elements $m,n \in M$, define the ideal $I_{m,n} \subset A$:
\[I_{m,n} = \eq(1 \mapsto m, 1 \mapsto n: A \to M),\]
and 
\[J_{m,n}=\eq(1 \mapsto m \otimes 1, 1 \mapsto n\otimes 1: B \to M \otimes_A B).\]
As $B$ is flat over $A$, we see that:
\[I_{m,n}B = I_{m,n} \otimes_A B = \eq(1 \mapsto m \otimes 1, 1 \mapsto n\otimes 1: B \to M\otimes_A B) = J_{m,n}.\]
We must show $J_{m,n} \neq B$ for $m \neq n$. If $m \neq n$, then $I_{m,n} \neq A$, so $I_{m,n}$ is contained in a maximal ideal $\ideal{m} \subset A$. Hence, 
$J_{m,n} = I_{m,n}B \subset \ideal{m}B$. By our assumption that $\text{Spec}(B) \to \text{Spec}(A)$, $\ideal{m}B \neq B$, and we therefore conclude
 that $J_{m,n} \neq B$.
\end{proof}
Hence, if $\phi\otimes_A B$ is injective, then by examining the diagram:
\[\begin{tikzcd} M \arrow[r, "\phi"] \arrow[hookrightarrow,d] & N\arrow[hookrightarrow,d]\\
    M \otimes_A B \arrow[r, "\phi \otimes_A B"] & N\otimes_A B 
\end{tikzcd}\]
we conclude that $\phi$ is injective.\\
\indent Let us now show that $\phi$ is surjective. For any $n \in N$, define the module $I_{n}$ as fitting into the pullback diagram:
\[
\begin{tikzcd}
I_n \arrow[r] \arrow[d]
  & M \arrow[d, "\phi"] \\
A \arrow[r, "1 \mapsto n"]
  & N
  \arrow[ul, phantom, "\lrcorner", very near start]
\end{tikzcd}
\]
and $J_{n\otimes 1}$ as fitting into the pullback diagram:
\[
\begin{tikzcd}
J_{n \otimes 1} \arrow[r] \arrow[d]
  & M\otimes_A B \arrow[d, "\phi \otimes_A B"] \\
B \arrow[r, "1 \mapsto n \otimes 1"]
  & N \otimes_A B
  \arrow[ul, phantom, "\lrcorner", very near start]
\end{tikzcd}
\]
By our assumption that $\phi \otimes_A B$ is an isomorphism, we see that $J_{n \otimes 1} = B$. Since $\phi$ is injective, 
$I_{n}$ is an ideal of $A$, and we must show $1 \in I_{n}$. As $A \to B$ is flat, 
\[I_nB = I_n \otimes_A B = J_{n\otimes1} = B.\]
By our assumption that $\text{Spec}(B) \to \text{Spec}(A)$ is surjective on maximal ideals, we see that $I_n$ is not contained in any maximal ideal of $A$, 
and therefore conclude $I_n = A$, so $1 \in I_n$ as desired.
\end{proof}
\begin{proposition}\label{prop:ff}
Let $f: A \to B$ is a faithfully flat map of semirings, then $\text{Spec}(B) \to \text{Spec}(A)$ is surjective.
\end{proposition}
\begin{proof}
We first show that $\text{Spec}(f)$ is surjective on maximal ideals. In particular, let $\ideal{m} \subset A$ be a maximal ideal and suppose $\ideal{m}B \subsetneq B$, 
then we may take a maximal ideal $\ideal{m}_B$ with $\ideal{m}B \subset \ideal{m}_B \subsetneq B$, and we see 
\[\text{Spec}(f)(\ideal{m}_B) = \ideal{m}\]
Hence, surjectivity on maximal ideals follows from the following claim:
\begin{claim}\label{claim:ff1}
Let $I \subsetneq A$ be a proper ideal. Then $IB=I\otimes_A B \neq B$.
\end{claim}
\begin{proof}[Proof of claim~\ref{claim:ff1}]
Suppose for the sake of contradiction that we may write $1 = \sum_{i=1}^{n} a_i \otimes b_i$ where $a_i \in I$ and $b_i \in B$, and let $\phi: A^{\oplus n} \to A$ denote the map sending $e_i \mapsto a_i$. Then we note that $\phi \otimes_A B$ is surjective. 
Because $A \to B$ is faithfully flat, this means $\phi$ is surjective, which is a contradiction.
\end{proof}
Let $\ideal{p} \subset A$ be any ideal of $A$. Consider the diagram:
\[\begin{tikzcd}
\text{Spec}(B) \arrow[r] & \text{Spec}(A) \\
\text{Spec}(B_{\ideal{p}}) \arrow[u] \arrow[r] & \text{Spec}(A_{\ideal{p}}) \arrow[u]
\end{tikzcd}\]
By what we just proved, the maximal ideal $\ideal{p}A_{\ideal{p}} \in \text{Spec}(A_{\ideal{p}})$ is in the image of the bottom-right arrow. Hence, $\ideal{p}$ is in the image of the top-right arrow, as desired.
\end{proof}
\begin{corollary}\label{cor:ff=fsurj}
A map $f: A \to B$ of semirings is faithfully flat if and only if it is flat and $\text{Spec}(f)$ is surjective.
\end{corollary}
Unsurprisingly now, we have going-down for flat ring maps.
\begin{theorem}[Going-down for semirings]\label{theorem:goingdown}
Let $f: A \to B$ be a flat map of semirings. Then the image of $\text{Spec}(f)$ is closed under generalizations.
\end{theorem}
\begin{proof}
Note that if $\varphi: R \to S$ is a map of \textit{local} semirings that is flat, it is automatically faithfully flat. Indeed, by supposition $\text{Spec}(\varphi)$ is surjective on maximal ideals, 
therefore by Proposition~\ref{prop:flatsurjmax}, is faithfully flat.\\
\indent Now let $\ideal{q} \subset B$ be any ideal mapping to an ideal $\ideal{p} \subset A$. Then the corresponding map $\text{Spec}(B_{\ideal{q}}) \to \text{Spec}(A_{\ideal{p}})$ is surjective, and therefore for any 
generalization $\ideal{p}' \subset \ideal{p} \subset A$ there is an ideal $\ideal{q}' \subset \ideal{q} \subset B$ such that $\text{Spec}(f)(\ideal{q}') = \ideal{p}'$
\end{proof}
We now have a proof of our main theorem.
\begin{theorem}\label{theorem:main}
A flat, finitely presented epimorphism of semirings $f: A \to B$ induces an isomorphism of $\text{Spec}(B)$ onto a quasi-compact open of $\text{Spec}(A)$ in the Zariski-topology.
\end{theorem}
\begin{proof}
Let $\ideal{p} \in \text{Spec}(A)$ be an prime ideal in the image of $\text{Spec}(f)$. Then $f \otimes_{A} A_{\ideal{p}}: A_{\ideal{p}} \to B \otimes_{A} A_{\ideal{p}}$ 
is a flat, finitely presented epimorphism such that $\text{Spec}(f \otimes_{A} A_{\ideal{p}})$ is surjective on maximal ideals. 
Hence, by Proposition~\ref{prop:flatsurjmax}, $f \otimes_{A} A_{\ideal{p}}$ is faithfully flat too, and is therefore an isomorphism.\\
\indent We next claim that there exists a $g \in A \setminus \ideal{p}$ such that $A_{g} \to B_g$ is an isomorphism, 
which is sufficient to conclude the proof of the theorem. Hence, it suffices to show:
\begin{claim}\label{claim:main1}
Let $I$ is a small filtered category, $F: I \to A\mathrm{-alg}_{\text{fp}}$ be a functor. 
Let $f: A \to B$ be a map in $A\mathrm{-alg}_{\text{fp}}$, and for each $i \in I$, let $f_i = f \otimes_A A_i$. 
If $\colim_{i \in I} f_i$ is an isomorphism, then there is an $i \in I$ such that $f_i$ is an isomorphism.
\end{claim}
\begin{proof} 
The proof is a standard result in category theory, but for the sake of convenience, we provide a proof.\\
\indent Let $B_i = B \otimes_A A_i$, $A_{\infty}=\colim_{i \in I} A_{\infty}, B_{\infty}=\colim_{i \in I}B_i$ and 
$f_{\infty}=\colim_{i \in I} f_i$, and let $g: B_{\infty} \to A_{\infty}$ be $A_{\infty}$-algebra inverse for $f$.\\
\indent By standard compactness arguments, there is an $i \in I$ and a map of $A_i$-algebras $g_i: B_i \to A_i$ such that
$g = g_i \otimes_{A_i} A_{\infty}$. Again by compactness, there is a $j \in I$ with $i \to j$ such that the composition
$g_i \otimes_{A_i} A_j \circ f_j: A_j \to A_j$ is the identity. By a symmetric argument, there is a $k \in I$ with $j \to k$ such that
$f_k \circ g_i \otimes_{A_i} A_k: B_k \to B_k$ is the identity, so $f_k$ is an isomorphism, as desired.
\end{proof}
\end{proof}
\begin{remark}[Images of finitely presented morphisms of semiring schemes are not constructible]\label{rem:badclosed}
One may ask if the epimorphism condition in Theorem~\ref{theorem:main} is necessary for the image of $\spec{f}$ to be open?
The usual method in classical commutative algebra to prove something like this rests on the fact that images of finitely presented
ring maps are constructible subsets (\cite[Tag 054k]{sp}), and then one uses going-down to conclude openness. However, consider the following map of semirings:
\[f: \bb{R}_{+}[x] \to \bb{R}_+, \; x \mapsto 1.\]
The map is finitely presented as it is cut-out by the equation $x=1$ in $\bb{R}_{+}[x]$. The spectrum of $\bb{R}_{+}$ has a unique prime ideal $\{0\} \subset \bb{R}_{+}$. 
The interesting fact about $f$ that happens to be true in the semiring setting is that $f^{-1}(\{0\}) = \{0\} \subset \bb{R}_{+}[x]$, so $\{0\}$ is the image of $\spec{f}$. 
But $\{0\}$ is not a constructible subset of $\spec{\bb{R}_{+}[x]}$. Indeed, if it were, then since $\{0\}$ is the generic point of $\spec{\bb{R}_{+}[x]}$, 
it would have to contain a non-empty open subset of $\spec{\bb{R}_{+}[x]}$\textemdash this is because $\spec{\bb{R}_{+}[x]}$ is a spectral space by Proposition~\ref{prop:spec} so 
we may appeal to \cite[Tag 0AAW]{sp}. 
Hence, $\{0\}$ would have to be an open subset of $\spec{\bb{R}_{+}[x]}$, which means that there is an $f \in \bb{R}_{+}[x]$ such that $D(f) = \{0\}$. 
But for every $a \in \bb{R}_{+}$, the ideal $(x+a) \subset \bb{R}_{+}[x]$ is prime, so since there is no polynomial $f$ that vanishes at $-a$ for all $a \in \bb{R}_{+}$, 
we conclude that $D(f) \neq \{0\}$, a contradiction.
\end{remark}
We will end this section with some further examples of how closed subschemes behave strangely in semiring algebraic geometry.
\begin{example}[Closed subschemes of semiring schemes behave unexpectedly]\label{exmp:badclosed}
Let $A=\bb{R}_{+}[x,y]/(x^2=x,y^2=y,1+xy=x+y)$. All elements of $A$ are of the form $a + bx + cy + dxy$ for $a,b,c,d \in \bb{R}_{+}$, 
and due to the relations, we may suppose at least one of the four coefficients is zero. The units are precisely the elements with $a \neq 0$ and $b=c=d=0$,
and there are no zero-divisors. We have the following prime ideals:
\begin{enumerate}
\item[(1)] The unique minimal prime $(0)$.
\item[(2)] The height one primes $(x)$, $(y)$.
\item[(3)] Height two primes of the form 
\[\ideal{p}=\{a + bx + cy + dxy| \; (a=0) \lor (d>0) \lor (a>0, b>0, c=d=0)\},\]
\[\ideal{q}=\{a + bx + cy + dxy| \; (a=0) \lor (d>0) \lor (a>0, c>0, b=d=0)\}.\]
\item[(4)] The unique maximal ideal $\ideal{m}$ consisting of the non-units.
\end{enumerate}
Therefore, $\spec{A}$ has Krull dimension $3$.\\
\indent A similar characterisation shows that $A[1/y]=\bb{R}_{+}[x]/(x=x^2)$ and $A[1/x]=\bb{R}_{+}[y]/(y=y^2)$ 
both have Krull dimension $2$. Note that $A \to A[1/x] \times A[1/y]$ is injective, and each $x$ and $y$ is integral over $A$, 
so $A \to A[1/x] \times A[1/y]$ is an integral extension of semirings. However, the map $\spec{A[1/x]} \coprod \spec{A[1/y]} \to \spec{A}$ 
is not surjective, in contrast to the usual going-up theorem in classical commutative algebra.\\
\indent In fact, for any $\ideal{p}$ of type $3$ above, there is no map of semirings $f:A \to B$ with $\ideal{p} = \eq(f,0:A \to B)$. 
In particular, both $xy$ and $1+xy$ are in $\ideal{p}$, so
\[1=1+f(xy)=f(1+xy)=0\]
in $B$, so $B$ is a trivial semiring, which would imply $\ideal{p}= A$, a contradiction.
\end{example}
\section{Flatness of quotient maps}\label{sect:flatquot}
This section is concerned with understanding when quotients $\bb{R}_{+}[x_1,...,x_n]/(f \sim g)$ are flat over $\bb{R}_{+}$. 
This is a surprisingly subtle question, even in the case of one variable polynomials.\\
\indent We begin with first an explicit construction of quotients of semirings.
\begin{definition}
Let $A$ be a semiring, $I$ a set, and for each $i \in I$, let $(f_i,g_i) \in A \times A$ be pairs of elements.
Let $\cal{R}(f_i\sim g_i)_{i\in I}$ be the smallest equivalence relation in $A\times A$ containing the elements $(f_i,g_i)$ that forms a subalgebra of $A \times A$.
Then the set $A/\cal{R}(f_i\sim g_i)_{i\in I}$ of equivalence classes has a natural semiring structure induced from $A$, 
and we call it the \textit{quotient} of $A$ by the relations $(f_i \sim g_i)_{i \in I}$.\\
\indent Let us show that $A/\cal{R}(f_i\sim g_i)_{i\in I}$ has the following universal property: for any semiring $B$ and any map $\phi: A \to B$ such 
that $\phi(f_i) = \phi(g_i)$ for all $i \in I$, there is a unique map $\overline{\phi}: A/\cal{R}(f_i\sim g_i)_{i\in I} \to B$ such that the diagram:
\[\begin{tikzcd} A \arrow[r, "\phi"] \arrow[d] & B \\
    A/\cal{R}(f_i\sim g_i)_{i\in I} \arrow[ur, "\overline{\phi}"]  & 
\end{tikzcd}\]
commutes. Let $\mathrm{diag}(B) \subset B \times B$ denote the trivial equivalence relation, clearly $\mathrm{diag}(B)$ is a sub-semiring. Consider the map 
$\phi \times \phi: A \times A \to B \times B$. Since $\phi(f_i) = \phi(g_i)$, we have $(f_i,g_i) \in (\phi \times \phi)^{-1}(\mathrm{diag}(B))$.
Moreover, $(\phi \times \phi)^{-1}(\mathrm{diag}(B))$ forms an equivalence relation, and is clearly a sub-semiring of $A \times A$,
so by minimality of $\cal{R}(f_i\sim g_i)_{i\in I}$, we have $\cal{R}(f_i\sim g_i)_{i\in I} \subset (\phi \times \phi)^{-1}(\mathrm{diag}(B))$. 
Hence, the map $\overline{\phi}$ taking $[a] \mapsto \phi(a)$ is well-defined, 
and is clearly unique.
\end{definition}
The set $\cal{R}(f_i \sim g_i)_{i \in I}$ is very hard to describe, and in my view,
is the primary reason for flatness and closed immersions in semiring algebraic geometry presenting many difficulties. In the case of one 
relation, we are able to be explicit.
\begin{proposition}\label{prop:quotientrels}
Let $f,g \in \bb{R}_{+}[x_1,...,x_n]$ be two elements and let $A=\bb{R}_{+}[x_1,...,x_n]/(f \sim g)$. The equivalence 
relation $\cal{R}(f \sim g) \subset \bb{R}_{+}[x_1,...,x_n] \times \bb{R}_{+}[x_1,...,x_n]$ 
generated by $(f \sim g)$ consists of pairs $(a_0 + b_0f + c_0g, a_t + b_tg + c_tf)$ for $a_0,b_0,c_0, a_t,b_t,c_t \in \bb{R}_{+}[x_1,...,x_n]$
such that there is a sequence of elements $(a_1,b_1,c_1),...,(a_{t-1},b_{t-1},c_{t-1})$ in $\bb{R}_{+}[x_1,...,x_n]$ with:
\[a_{i} + b_{i}g + c_{i}f=a_{i+1} + b_{i+1}f + c_{i+1}g \]
as elements in $\bb{R}_{+}[x_1,...,x_n]$ for all $0 \le i \le t-1$.
\end{proposition}
\begin{proof}
Let $\cal{R}'$ be the proposed relation. First, we note that $\cal{R}'$ is by design transitive, reflexive and symmetric, and hence an equivalence relation.
On the other hand, by reflexivity and symmetry, elements of the form $(a + bf + cg, a + bg + cf)$ are in $\cal{R}(f \sim g)$ for any $a,b,c \in A$.
Moreover, by transitivity, if there is a sequence of elements $(a_1,b_1,c_1),...,(a_{t-1},b_{t-1},c_{t-1})$ in $A$ with:
\[a_{i} + b_{i}g + c_{i}f= a_{i+1} + b_{i+1}f + c_{i+1}g\]
as elements in $\bb{R}_{+}[x_1,...,x_n]$ for all $0 \le i \le t-1$, then $(a_0 + b_0f + c_0g, a_t + b_tg + c_tf) \in \cal{R}(f\sim g)$. Therefore, $\cal{R}' \subset \cal{R}(f \sim g)$, and to finish,
 it suffices to show that $\cal{R}'$ is a sub-semiring of $A \times A$.\\
\indent For any $d \in \bb{R}_{+}[x_1,...,x_n]$ and $(h,h') \in \cal{R}'$, then $(dh,dh')\in \cal{R}'$, because the triples $\{(a_i,b_i,c_i)\}_{i=0}^t$
realizing the equality $(h,h')\in \cal{R}'$ may be multiplied by $d$ to realize $(dh,dh') \in \cal{R}'$. Therefore, for any two elements
$(h_1,h'_1), (h_2,h'_2) \in \cal{R}'$, we have $(h_1h_2,h_1h'_2) \in \cal{R}'$ and $(h'_1h_2,h'_1h'_2) \in \cal{R}'$, 
so $(h_1h_2,h'_1h'_2) \in \cal{R}'$
by transitivity. 
\end{proof}
\begin{corollary}\label{cor:univpointless}
The semiring $A:=\bb{R}_{+}[x_1,y_1]/(1+x_1^2+y_1^2 \sim 2x_1y_1)$ is not flat over $\bb{R}_{+}$.
\end{corollary}
\begin{proof}
Let $\cal{R}$ be the equivalence relation generated by $1+x_1^2 + y_1^2 \sim 2x_1y_1$. Using Proposition~\ref{prop:quotientrels}, we may observe that if $(a,b) \in \cal{R}$, then
$\mathrm{deg}(a)=\mathrm{deg}(b)$. Note that this would more generally hold for any relation of the form $p \sim q$ where $p$ and $q$ are polynomials of the same degree. 
Therefore, by matching top-degrees, we see that if $(a,b) \in \cal{R}$ and $a'$ and $b'$ are the top-degree terms of $a$ and $b$ respectively,
then $(a',b') \in \cal{R}'$, where $\cal{R}'$ is the equivalence relation generated by $x_1^2+y_1^2 \sim 2x_1y_1$.\\
\indent Since $\bb{R}_{+}$ is cancellative, we know by Proposition~\ref{prop:flatpositivity} that if $A$ is flat, then it is also cancellative.
We have two equations:
\[2x_1 + 2x_1^3 + 2x_1y_1^2 \sim 4x_1^2y_1\]
\[y_1 + x_1^2y_1 +  y_1^3 \sim 2x_1y_1^2,\]
so by substitution, cancellativity and matching top-degrees, we find that $(2x_1^3 + y_1^3,3x_1^2y_1)$ is in the degree $3$ part of $\cal{R}'$, which is generated by the relations:
\[x_1^3 + x_1y_1^2 \sim 2x_1^2y_1\]
\[x_1^2y_1 + y_1^3 \sim 2x_1y_1^2.\]
For arbitrary elements $h=a_1x_1^3 + a_2x_1^2y_1 + a_3x_1y_1^2 + a_4y_1^3$ and $k=b_1x_1^3 + b_2x_1^2y_1 + b_3x_1y_1^2 + b_4y_1^3$, if 
$(h,k)$ is in the degree $3$ part of $\cal{R}'$, then by Proposition~\ref{prop:quotientrels}, if $h$ and $k$ are not equal as polynomials, 
then because both degree $3$ generators for $\cal{R}'$ contain non-zero coefficients for $x_1y_1^2$ and $x_1^2y_1$ (on both sides), we see that
$a_2 + a_3$ and $b_2 + b_3$ are both strictly positive. Therefore, the pair $(2x_1^3 + y_1^3,3x_1^2y_1)$ cannot be in the degree $3$ part of $\cal{R}'$, so we have a contradiction, 
and $A$ is not flat over $\bb{R}_{+}$.
\end{proof}
\begin{remark}
Corollary~\ref{cor:univpointless} shows that the obvious choice for a counterexample to Question~\ref{mainquestion} does not work. 
A natural next step would be to prove that the higher dimensional analogue of the relation $1+x_1^2+y_1^2 \sim 2x_1y_1$, which is 
$1+f \sim g$ where $f=\sum_{i \in [n]} x_i^2 + y_i^2$ and $g=\sum_{i \in [n]} 2x_iy_i$, does not generate a $\bb{R}_{+}$-flat quotient
$A=\bb{R}_{+}[x_1,y_1,...,x_n,y_n]/(1+f \sim g)$ either. We did not pursue this question seriously, but we outline a possible approach mimicking the proof for $n=1$.\\
\indent Indeed, the key idea was to use cancellativity to find a relation that cannot be generated by $f \sim g$ alone. We could consider the following two equations:
\[g+(f-y_n^2)g + y_n^2g \sim g^2\]
\[y_n^2+ y_n^2f \sim y_n^2g,\]
so by substitution, cancellativity and matching top-degrees, we find that $(h,k)$ is in the degree $4$ part of $f \sim g$, where 
\[h=(f-y_n^2)(g+y_n^2)- x_n^2y_n^2 + y_n^4 \]
\[k=g^2 - x_n^2y_n^2.\]
Let $A_n \subset A$ denote the set of homogenous degree $n$ monomials. The degree $4$ part of $\cal{R}'$ is the transitive and symmetric closure of the relations
generated by the equations
\[af \sim ag, a \in A_2.\]
These equations however do not exhibit the same kind of invariant as the degree $3$ relations in the $n=1$ case, so it is not clear how to show that $(h,k)$ is not in the degree $4$ part of $\cal{R}'$.
\end{remark}
The fundamental tool to analyze flatness is the equational criterion for flatness, which holds in the semiring setting as well (see \cite[3.5]{borgersemi}).
\begin{theorem}[Equational criterion for flatness]\label{theorem:equationalflatness}
Let $A$ be a semiring, and let $M$ be an $A$-module. Then $M$ is flat if and only if for any finite collection of elements $m_1,...,m_k \in M$ satisfying a relation:
\[\sum_{i\in [k]} a_{i}m_i = \sum_{i\in [k]} b_{i}m_i,\]
with $a_{i}, b_{i} \in A$, there is a finite set of elements $\{w_l\}_{l\in [t]}, \subset M$, and for each $i\in[k], l \in [t]$,
there are elements $c_{il} \in A$ such that:
\[m_i = \sum_{l\in [t]} c_{il}w_l = \sum_{l\in [t]} d_{il}w_l\]
\[\sum_{i\in [k]} a_{i}c_{il} = \sum_{i\in [k]} b_{i}c_{il}, \; \forall l \in [t].\]
Consequently, $M$ is flat if and only if it can be written as a filtered colimit of free $A$-modules.
\end{theorem}
Here are a few consequences.
\begin{corollary}\label{cor:flatquotient}
Let $A$ be a semiring, $B=A[x_1,...,x_n]/(p_i \sim q_i)_{i \in [k]}$ be a quotient of a polynomial semiring over $A$. 
If $B$ is flat, then the following are true:
\begin{enumerate}
\item[(i)] Without changing the choice of surjection $A[x_1,...,x_n] \to B$, we may assume $p_i$ is a monomial for all $i \in [k]$.
\item[(ii)] Now let $A = \bb{N}$, and suppose that no monomial goes to $0$ in $B$. 
Using (i), for each $i$ write $p_i=x^{E_i}$ and $q_i = \sum_{j\in [m_i]} a_i x^{F_{ij}}$ for $E_i, F_{ij} \in \bb{N}^n$. 
Then for each $i \in [k]$, there do not exist $\{a_j\}_{j \in [m_i]} \subset \bb{N}$ such that
\[\sum_{j \in [m_i]} a_j (E_i-F_{ij}) = 0,\]
unless $a_j = 0$ for all $j \in [m_i]$. In particular, $E_i$ is not in the convex hull of $\{F_{ij}\}_{j \in [m_i]}$.
\end{enumerate}
\end{corollary}
We first prove a lemma.
\begin{lemma}[Boosting polynomial relations using flatness]\label{lemma:flatmonomial}
Let $A$ be a semiring, and let $\phi: A[x_1,...,x_n] \to B$ a surjection of semirings. Suppose $B$ is flat over $A$. 
Then for any two elements $f,g \in A[x_1,...,x_n]$ such that $\phi(f)=\phi(g)$, there is a finite set of elements $t_i,s_i \in A[x_1,...,x_n]$ with $\phi(t_i)=\phi(s_i)$ for all $i$,
such that each $t_i$ is a monomial, and the relation $f \sim g$ is implied by the relations $t_i \sim s_i$ for all $i$.
\end{lemma}
\begin{proof}
Write $f= \sum_{i\in [m]} a_i x^{E_i}$ and $g = \sum_{j\in [l]} b_j x^{F_j}$ for $a_i,b_j \in A$ and $E_i, F_j \in \bb{N}^n$. Since $f$ and $g$ have the same image in $B$, 
by Theorem~\ref{theorem:equationalflatness}, 
there are elements $w_1,...,w_t \in A[x_1,...,x_n]$, elements $\{c_{ik}\}_{i \in [m], k \in [t]}$ and $\{d_{jk}\}_{j \in [l], k \in [t]}$ in $A$ such that:
\[\phi(x^{E_i}) = \phi\left(\sum_{k \in [t]} c_{ik} w_k\right), \; \phi(x^{F_j}) = \phi\left(\sum_{k \in [t]} d_{jk} w_k\right),\]
and for each $k \in [t]$,
\[\sum_{i\in [m]} a_i c_{ik} = \sum_{j\in [l]} b_j d_{jk}.\]
By setting $t_i=x^{E_i}$ for $i \in [m]$ and $t_{m+j} = x^{F_j}$ for $j \in [l]$, and $s_i = \sum_{k \in [t]} c_{ik} w_k$ for $i \in [m]$ and
$s_{m+j} = \sum_{k \in [t]} d_{jk} w_k$ for $j \in [l]$, we conclude the proof.
\end{proof}
\begin{proof}[Proof of Corollary~\ref{cor:flatquotient}]
For (i), we may apply Lemma~\ref{lemma:flatmonomial} to each relation $p_i \sim q_i$ to obtain the desired result.\\
\indent For (ii), recall that $B$ is an integral commutative monoid (see e.g. Proposition~\ref{prop:flatpositivity}), and 
can therefore be endowed with a partial ordering by Proposition~\ref{prop:poset}. We will endow the space $\bb{N}^n$ with the following ordering. 
For two $E,F \in \bb{N}^n$, we say $E \le F$ if and only if $x^E \le x^F$ in $B$. It's clear that $(\bb{N}^n, \le)$ defines
a partial ordering, and since $B$ is integral as a module over itself, $\le$ is compatible with addition.\\
\indent Now suppose for the sake of contradiction that for some $i \in [k]$, there are numbers $\{a_j\}_{j \in [m_i]} \subset \bb{N}$ such that:
\[\sum_{j\in [m_i]} a_j (E_i-F_{ij}) = 0.\]
This implies that:
\[\left(\sum_{j\in [m_i]} a_j\right) E_i = \sum_{j\in [m_i]} a_j F_{ij}.\]
On the other hand, the equation $p_i=q_i$ implies that $E_i > F_{ij}$ for all $1 \le j \le m_i$. Therefore, we have:
\[\left(\sum_{j\in [m_i]} a_j\right) E_i > \sum_{j\in [m_i]} a_j F_{ij} = \left(\sum_{j\in [m_i]} a_j\right) E_i,\]
a contradiction unless $\sum_{j\in [m_i]} a_j=0 \implies a_j=0$ for all $j \in [m_i]$.
\end{proof}
\begin{corollary}
Let $A=\bb{N}[x_1,...,x_n]/(x^{E}\sim \sum_{j \in [m]} a_j x^{F_j})$ for $E,F_j \in \bb{N}^n$ and $a_j \in \bb{N}$. 
If $A$ is flat over $\bb{N}$, then
$A$ has an $\bb{R}_{+}$-point.
\end{corollary}
\begin{proof}
By Corollary~\ref{cor:flatquotient} (ii) and the Hahn-Banach theorem, there is a linear functional $\lambda: \bb{R}^n \to \bb{R}$ such that
\[\lambda(E) > \lambda(F_j), \; \forall j \in [m].\]
By perturbation we may assume $\lambda: \bb{Q}^n \to \bb{Q}$, and by scaling we may assume $\lambda: \bb{Z}^n \to \bb{Z}$. Define the semiring map
\[\phi: \bb{N}[x_1,...,x_n] \to \bb{N}[t^{\pm 1}]\]
induced by sending $x_i \mapsto t^{\lambda(e_i)}$ where $\{e_i\}_{i=1}^n$ is the standard basis of $\bb{Z}^n$. Let
\[A'=\bb{N}[t^{\pm1}]/(t^{\lambda(E)} \sim \sum_{j \in [m]} a_j t^{\lambda(F_j)}).\]
Then $\phi$ induces a map $A \to A'$, and since $\lambda(E) > \lambda(F_j)$ for all $j \in [m]$, the equation
\[t^{\lambda(E)} = \sum_{j \in [m]} a_j t^{\lambda(F_j)}\]
has at least one solution in $\bb{R}_{+}$. Therefore, $A'$ has an $\bb{R}_{+}$-point, and hence so does $A$.
\end{proof}
